\documentclass[10pt]{article}
 \usepackage[margin=1in]{geometry} 
\usepackage{amsmath,amsthm,amssymb,amsfonts,color, titling}
 \usepackage{listings}

\setlength{\droptitle}{-20mm} 

\usepackage[colorlinks=true,linkcolor=red,urlcolor=blue]{hyperref}

\title{Pre-class assignment \#21}
\author{PHY-905-003\\Computational Astrophysics and
  Astrostatistics\\Spring 2026}
 \date{} % leave blank to have no date

\begin{document}
 
\maketitle

\vspace{-10mm}

\noindent
\textbf{This assignment is due the evening of Monday March 30, 2026.}
 Turn in all materials via GitHub.

\vspace{5mm}

\noindent
\textbf{Reading:}

\begin{enumerate}

\item Chapters 1 and 2 of
  \href{https://artowen.su.domains/mc/}{\textit{Monte Carlo
      Theory, Methods, and Examples}} by Art Owen.  \href{https://artowen.su.domains/mc/Ch-intro.pdf}{Direct
    link to chapters.}  (Chapters included in this repository as PDF.)

\item The
  \href{https://en.wikipedia.org/wiki/Monte\_Carlo\_integration}{Wikipedia
    page on Monte Carlo Integration}

\item The
  \href{https://en.wikipedia.org/wiki/Pseudorandom_number_generator}{Wikipedia
  page on Pseudorandom number generators} (which is what most
computers use to generate ``random'' numbers).  Python and most other
programming languages use a generator known as the \href{https://en.wikipedia.org/wiki/Mersenne_Twister}{Mersenne
  Twister}, which is good enough for most non-cryptographic applications.

\end{enumerate}

\noindent
\textbf{Your assignment:}

\begin{enumerate}

\item Work through the tasks in the included Jupyter notebook!

\item In the file \texttt{ANSWERS.md}, write down any questions that
you have about the material you read or the work you did in the
Jupyter notebook discussed above, any points that are not clear, or
anything you'd like to know more about.  Aim for at least 3
questions/unclear points/etc.

\end{enumerate}

\noindent \textbf{Handing it in:} Include your modified code and
your answers to the questions (in the file \texttt{ANSWERS.md})
in your assignment.

\end{document}
